\documentclass[12pt]{ximera}
\typeout{Start loading xmPreamble.tex}%

% Add here extra macro's that are loaded automatically by all documents of claas 'ximera' or 'xourse' in this repo

%%
%%  Example:
%%
% \newcommand{\R}{\mathbb{R}}
\usepackage{tikz}
\usetikzlibrary{positioning, arrows.meta}
\usepackage{multicol}
\usepackage{enumitem}
\usepackage{caption}
\usepackage{amsmath}
\usepackage{amsthm}
\usepackage{fontenc}

%% ---------------------------------------------------------------
%% Shared worksheet macros.
%%
%% These came from the Summer 2026 worksheet set, where each worksheet carried its
%% own copy. They have to live here instead: a xourse inlines every activity into a
%% single document, so duplicate \newcommand definitions across activities collide,
%% and the xourse gobbles \input so a per-topic macro file never loads.
%%
%% They use \providecommand, NOT \newcommand. ximera.cls loads this file automatically,
%% and every activity also carries an explicit \typeout{Start loading xmPreamble.tex}%

% Add here extra macro's that are loaded automatically by all documents of claas 'ximera' or 'xourse' in this repo

%%
%%  Example:
%%
% \newcommand{\R}{\mathbb{R}}
\usepackage{tikz}
\usetikzlibrary{positioning, arrows.meta}
\usepackage{multicol}
\usepackage{enumitem}
\usepackage{caption}
\usepackage{amsmath}
\usepackage{amsthm}
\usepackage{fontenc}

%% ---------------------------------------------------------------
%% Shared worksheet macros.
%%
%% These came from the Summer 2026 worksheet set, where each worksheet carried its
%% own copy. They have to live here instead: a xourse inlines every activity into a
%% single document, so duplicate \newcommand definitions across activities collide,
%% and the xourse gobbles \input so a per-topic macro file never loads.
%%
%% They use \providecommand, NOT \newcommand. ximera.cls loads this file automatically,
%% and every activity also carries an explicit \typeout{Start loading xmPreamble.tex}%

% Add here extra macro's that are loaded automatically by all documents of claas 'ximera' or 'xourse' in this repo

%%
%%  Example:
%%
% \newcommand{\R}{\mathbb{R}}
\usepackage{tikz}
\usetikzlibrary{positioning, arrows.meta}
\usepackage{multicol}
\usepackage{enumitem}
\usepackage{caption}
\usepackage{amsmath}
\usepackage{amsthm}
\usepackage{fontenc}

%% ---------------------------------------------------------------
%% Shared worksheet macros.
%%
%% These came from the Summer 2026 worksheet set, where each worksheet carried its
%% own copy. They have to live here instead: a xourse inlines every activity into a
%% single document, so duplicate \newcommand definitions across activities collide,
%% and the xourse gobbles \input so a per-topic macro file never loads.
%%
%% They use \providecommand, NOT \newcommand. ximera.cls loads this file automatically,
%% and every activity also carries an explicit \input{../xmPreamble}, so this file is
%% read twice. That was harmless while it held only \usepackage lines (LaTeX skips a
%% package it has already loaded) but a second \newcommand is a hard error.
%%
%%   \blank[width]        fill-in-the-blank rule (default 2.5cm)
%%   \showwork            "show and explain your work" reminder
%%   \showworkline        ... plus which schedule line was used
%%   \showworkyear        ... plus which year's schedule was used
%%   \schedhead{...}      centred bold caption above a tiered-rate table
%%   \samesquares[scale]{left}{right}   two equal-area squares, labelled
%%   \samecubes[scale]{left}{right}     two equal-volume cubes, labelled
%% ---------------------------------------------------------------

\providecommand{\blank}[1][2.5cm]{\underline{\hspace{#1}}}
\providecommand{\showwork}{\textbf{REMEMBER TO SHOW AND EXPLAIN YOUR WORK.}}
\providecommand{\showworkline}{\textbf{REMEMBER TO SHOW AND EXPLAIN YOUR WORK.
  Explanation should include how and why you know which line of the
  schedule was used to calculate this amount.}}
\providecommand{\showworkyear}{\begin{sloppypar}\textbf{REMEMBER TO SHOW AND
  EXPLAIN YOUR WORK. Explanation should include how and why you know which
  line of the year's tax schedule was used to calculate this tax
  amount.}\end{sloppypar}}
\providecommand{\schedhead}[1]{\begin{center}\textbf{#1}\end{center}}

\providecommand{\samesquares}[3][1]{%
\begin{center}
{\footnotesize These squares are the \textbf{same size}.}

\vspace{1.5mm}
\begin{tikzpicture}[scale=#1,every node/.style={scale=#1}]
  \draw (0,0) rectangle (2.4,2.4);
  \node[anchor=north west] at (0.15,2.25) {Area:};
  \node[anchor=east]  at (-0.15,1.2) {#2};
  \node[anchor=north] at (1.2,-0.15) {#2};
  \begin{scope}[xshift=4.6cm]
    \draw (0,0) rectangle (2.4,2.4);
    \node[anchor=north west] at (0.15,2.25) {Area:};
    \node[anchor=east]  at (-0.15,1.2) {#3};
    \node[anchor=north] at (1.2,-0.15) {#3};
  \end{scope}
\end{tikzpicture}
\end{center}}

\providecommand{\samecubes}[3][1]{%
\begin{center}
{\footnotesize These cubes are the \textbf{same size}.}

\vspace{1.5mm}
\begin{tikzpicture}[scale=#1,every node/.style={scale=#1}]
  \foreach \dx in {0,5.2} {
    \begin{scope}[xshift=\dx cm]
      \draw (0,0) rectangle (2.0,2.0);
      \draw (0,2.0) -- (0.65,2.6) -- (2.65,2.6) -- (2.0,2.0);
      \fill[black!12] (2.0,0) -- (2.65,0.6) -- (2.65,2.6) -- (2.0,2.0) -- cycle;
      \draw (2.0,0) -- (2.65,0.6) -- (2.65,2.6) -- (2.0,2.0);
      \node[anchor=north west] at (0.1,1.9) {Volume:};
    \end{scope}
  }
  \node[anchor=east]  at (-0.15,1.0) {#2};
  \node[anchor=north] at (1.0,-0.15) {#2};
  \node[anchor=west]  at (2.25,0.4) {#2};
  \node[anchor=east]  at (5.05,1.0) {#3};
  \node[anchor=north] at (6.2,-0.15) {#3};
  \node[anchor=west]  at (7.45,0.4) {#3};
\end{tikzpicture}
\end{center}}

%% NOTE: do NOT add \usepackage to individual activity .tex files.
%% When a xourse inlines an activity it gobbles \documentclass and \input, but a bare
%% \usepackage still executes -- after \begin{document} -- and errors out.
%% All shared packages and macros belong here.
%%
%% ximera.cls already loads: amsmath, amssymb, amsthm, cancel, enumitem, graphicx,
%% hyperref, listings, pgfplots, tikz, url, xcolor. No need to re-load those.
, so this file is
%% read twice. That was harmless while it held only \usepackage lines (LaTeX skips a
%% package it has already loaded) but a second \newcommand is a hard error.
%%
%%   \blank[width]        fill-in-the-blank rule (default 2.5cm)
%%   \showwork            "show and explain your work" reminder
%%   \showworkline        ... plus which schedule line was used
%%   \showworkyear        ... plus which year's schedule was used
%%   \schedhead{...}      centred bold caption above a tiered-rate table
%%   \samesquares[scale]{left}{right}   two equal-area squares, labelled
%%   \samecubes[scale]{left}{right}     two equal-volume cubes, labelled
%% ---------------------------------------------------------------

\providecommand{\blank}[1][2.5cm]{\underline{\hspace{#1}}}
\providecommand{\showwork}{\textbf{REMEMBER TO SHOW AND EXPLAIN YOUR WORK.}}
\providecommand{\showworkline}{\textbf{REMEMBER TO SHOW AND EXPLAIN YOUR WORK.
  Explanation should include how and why you know which line of the
  schedule was used to calculate this amount.}}
\providecommand{\showworkyear}{\begin{sloppypar}\textbf{REMEMBER TO SHOW AND
  EXPLAIN YOUR WORK. Explanation should include how and why you know which
  line of the year's tax schedule was used to calculate this tax
  amount.}\end{sloppypar}}
\providecommand{\schedhead}[1]{\begin{center}\textbf{#1}\end{center}}

\providecommand{\samesquares}[3][1]{%
\begin{center}
{\footnotesize These squares are the \textbf{same size}.}

\vspace{1.5mm}
\begin{tikzpicture}[scale=#1,every node/.style={scale=#1}]
  \draw (0,0) rectangle (2.4,2.4);
  \node[anchor=north west] at (0.15,2.25) {Area:};
  \node[anchor=east]  at (-0.15,1.2) {#2};
  \node[anchor=north] at (1.2,-0.15) {#2};
  \begin{scope}[xshift=4.6cm]
    \draw (0,0) rectangle (2.4,2.4);
    \node[anchor=north west] at (0.15,2.25) {Area:};
    \node[anchor=east]  at (-0.15,1.2) {#3};
    \node[anchor=north] at (1.2,-0.15) {#3};
  \end{scope}
\end{tikzpicture}
\end{center}}

\providecommand{\samecubes}[3][1]{%
\begin{center}
{\footnotesize These cubes are the \textbf{same size}.}

\vspace{1.5mm}
\begin{tikzpicture}[scale=#1,every node/.style={scale=#1}]
  \foreach \dx in {0,5.2} {
    \begin{scope}[xshift=\dx cm]
      \draw (0,0) rectangle (2.0,2.0);
      \draw (0,2.0) -- (0.65,2.6) -- (2.65,2.6) -- (2.0,2.0);
      \fill[black!12] (2.0,0) -- (2.65,0.6) -- (2.65,2.6) -- (2.0,2.0) -- cycle;
      \draw (2.0,0) -- (2.65,0.6) -- (2.65,2.6) -- (2.0,2.0);
      \node[anchor=north west] at (0.1,1.9) {Volume:};
    \end{scope}
  }
  \node[anchor=east]  at (-0.15,1.0) {#2};
  \node[anchor=north] at (1.0,-0.15) {#2};
  \node[anchor=west]  at (2.25,0.4) {#2};
  \node[anchor=east]  at (5.05,1.0) {#3};
  \node[anchor=north] at (6.2,-0.15) {#3};
  \node[anchor=west]  at (7.45,0.4) {#3};
\end{tikzpicture}
\end{center}}

%% NOTE: do NOT add \usepackage to individual activity .tex files.
%% When a xourse inlines an activity it gobbles \documentclass and \input, but a bare
%% \usepackage still executes -- after \begin{document} -- and errors out.
%% All shared packages and macros belong here.
%%
%% ximera.cls already loads: amsmath, amssymb, amsthm, cancel, enumitem, graphicx,
%% hyperref, listings, pgfplots, tikz, url, xcolor. No need to re-load those.
, so this file is
%% read twice. That was harmless while it held only \usepackage lines (LaTeX skips a
%% package it has already loaded) but a second \newcommand is a hard error.
%%
%%   \blank[width]        fill-in-the-blank rule (default 2.5cm)
%%   \showwork            "show and explain your work" reminder
%%   \showworkline        ... plus which schedule line was used
%%   \showworkyear        ... plus which year's schedule was used
%%   \schedhead{...}      centred bold caption above a tiered-rate table
%%   \samesquares[scale]{left}{right}   two equal-area squares, labelled
%%   \samecubes[scale]{left}{right}     two equal-volume cubes, labelled
%% ---------------------------------------------------------------

\providecommand{\blank}[1][2.5cm]{\underline{\hspace{#1}}}
\providecommand{\showwork}{\textbf{REMEMBER TO SHOW AND EXPLAIN YOUR WORK.}}
\providecommand{\showworkline}{\textbf{REMEMBER TO SHOW AND EXPLAIN YOUR WORK.
  Explanation should include how and why you know which line of the
  schedule was used to calculate this amount.}}
\providecommand{\showworkyear}{\begin{sloppypar}\textbf{REMEMBER TO SHOW AND
  EXPLAIN YOUR WORK. Explanation should include how and why you know which
  line of the year's tax schedule was used to calculate this tax
  amount.}\end{sloppypar}}
\providecommand{\schedhead}[1]{\begin{center}\textbf{#1}\end{center}}

\providecommand{\samesquares}[3][1]{%
\begin{center}
{\footnotesize These squares are the \textbf{same size}.}

\vspace{1.5mm}
\begin{tikzpicture}[scale=#1,every node/.style={scale=#1}]
  \draw (0,0) rectangle (2.4,2.4);
  \node[anchor=north west] at (0.15,2.25) {Area:};
  \node[anchor=east]  at (-0.15,1.2) {#2};
  \node[anchor=north] at (1.2,-0.15) {#2};
  \begin{scope}[xshift=4.6cm]
    \draw (0,0) rectangle (2.4,2.4);
    \node[anchor=north west] at (0.15,2.25) {Area:};
    \node[anchor=east]  at (-0.15,1.2) {#3};
    \node[anchor=north] at (1.2,-0.15) {#3};
  \end{scope}
\end{tikzpicture}
\end{center}}

\providecommand{\samecubes}[3][1]{%
\begin{center}
{\footnotesize These cubes are the \textbf{same size}.}

\vspace{1.5mm}
\begin{tikzpicture}[scale=#1,every node/.style={scale=#1}]
  \foreach \dx in {0,5.2} {
    \begin{scope}[xshift=\dx cm]
      \draw (0,0) rectangle (2.0,2.0);
      \draw (0,2.0) -- (0.65,2.6) -- (2.65,2.6) -- (2.0,2.0);
      \fill[black!12] (2.0,0) -- (2.65,0.6) -- (2.65,2.6) -- (2.0,2.0) -- cycle;
      \draw (2.0,0) -- (2.65,0.6) -- (2.65,2.6) -- (2.0,2.0);
      \node[anchor=north west] at (0.1,1.9) {Volume:};
    \end{scope}
  }
  \node[anchor=east]  at (-0.15,1.0) {#2};
  \node[anchor=north] at (1.0,-0.15) {#2};
  \node[anchor=west]  at (2.25,0.4) {#2};
  \node[anchor=east]  at (5.05,1.0) {#3};
  \node[anchor=north] at (6.2,-0.15) {#3};
  \node[anchor=west]  at (7.45,0.4) {#3};
\end{tikzpicture}
\end{center}}

%% NOTE: do NOT add \usepackage to individual activity .tex files.
%% When a xourse inlines an activity it gobbles \documentclass and \input, but a bare
%% \usepackage still executes -- after \begin{document} -- and errors out.
%% All shared packages and macros belong here.
%%
%% ximera.cls already loads: amsmath, amssymb, amsthm, cancel, enumitem, graphicx,
%% hyperref, listings, pgfplots, tikz, url, xcolor. No need to re-load those.

\title{Percentages: Reading ``Per'' and ``Cent''}
\begin{document}
\begin{abstract}
Building percentages up from the words themselves --- which English words signal which
operation, what ``cent'' means, and why that makes $13\%$ equal to $\tfrac{13}{100}$.
Then translating a ``portion is percent of whole'' sentence into an equation and
solving it.
\end{abstract}
\maketitle

\section*{Breaking Down the Word}

Percentages are everywhere! They are used to describe discounts, markups, commissions,
statistics, information, change, and on and on. They are an important topic and so we
want to make sure that we really understand them. Let's start by just breaking down the
word ``percent'': ``per'' and ``cent.''

There are some specific words that often translate into certain mathematical
operations. For example, if a question asks ``how many apples and oranges are there?''
what operation are you going to do with the number of apples and oranges?

\begin{problem}
``and'' often translates to:
\begin{multipleChoice}
\choice[correct]{addition}
\choice{subtraction}
\choice{multiplication}
\choice{division}
\end{multipleChoice}

``difference'' often translates to:
\begin{multipleChoice}
\choice{addition}
\choice[correct]{subtraction}
\choice{multiplication}
\choice{division}
\end{multipleChoice}

``of'' often translates to:
\begin{multipleChoice}
\choice{addition}
\choice{subtraction}
\choice[correct]{multiplication}
\choice{division}
\end{multipleChoice}

``per'' often translates to:
\begin{multipleChoice}
\choice{addition}
\choice{subtraction}
\choice{multiplication}
\choice[correct]{division}
\end{multipleChoice}

``is'' often translates to:
\begin{multipleChoice}
\choice[correct]{equals}
\choice{addition}
\choice{multiplication}
\choice{division}
\end{multipleChoice}

\begin{solution}
``and'' signals addition, ``difference'' signals subtraction, ``of'' signals
multiplication, ``per'' signals division, and ``is'' signals an equals sign. These five
translations are enough to turn most percentage sentences directly into equations.
\end{solution}
\end{problem}

\begin{problem}
Now ``cent.'' Where have we seen this word before? How many cents are in a dollar? How
many years are in a century? How many centimeters are in a meter?

Wherever you see the word ``cent,'' it is likely representing the number
$\answer{100}$.

\begin{solution}
100 cents in a dollar, 100 years in a century, 100 centimeters in a meter. ``Cent''
means 100.
\end{solution}
\end{problem}

\begin{problem}
Putting the two halves together: ``per'' means division and ``cent'' means 100, so
``percent'' can be interpreted as \emph{divided by 100}. This is why
\[
13\% \;=\; \frac{\answer{13}}{\answer{100}} \;=\; 0.13
\]

\begin{solution}
``Percent'' is literally ``per hundred,'' that is, divided by 100. So $13\%$ is
$\tfrac{13}{100}$, which as a decimal is $0.13$.
\end{solution}
\end{problem}

\section*{Turning a Sentence into an Equation}

This is also why ``$13\%$ of 200 is 26.'' Translate each piece:

\begin{center}
$13\% \;\to\; 0.13$ \qquad $\text{of} \;\to\; \times$ \qquad $200 \;\to\; 200$
\qquad $\text{is} \;\to\; =$ \qquad $26 \;\to\; 26$
\end{center}

Making ``$13\%$ of 200 is 26'' become the equation $0.13 \times 200 = 26$.

\begin{problem}
After a garage sale, you see that $35\%$ of what was sold was old records. If 14
records were sold, how many total items were sold at the sale?

To answer this question, we should take the values we know and write them in the
phrasing \emph{(portion) is (\%) of (whole)}. Filling in what we know:

$\answer{14}$ is $\answer{35}\%$ of (unknown).

\begin{solution}
14 records is the portion, $35\%$ is the percent, and the total number of items sold is
the whole --- which is what we are missing.
\end{solution}
\end{problem}

\begin{problem}
Now write that \emph{(portion) is (\%) of (whole)} sentence as a mathematical equation,
putting an $x$ in for the unknown. Then solve for $x$ and write a sentence about what
this value means in terms of the original question.

There were $\answer{40}$ total items sold at the sale.

\begin{freeResponse}
\end{freeResponse}

\begin{solution}
Translating word by word, ``14 is $35\%$ of $x$'' becomes
\[
14 = 0.35 \cdot x.
\]
Solving,
\[
x = \frac{14}{0.35} = 40.
\]
So 40 items were sold at the garage sale in total. Checking against the context: $35\%$
of 40 is $0.35 \cdot 40 = 14$, which matches the 14 records that were sold.
\end{solution}
\end{problem}

\end{document}
